% Part: computability
% Chapter: recursive-functions
% Section: general-recursive-functions

\documentclass[../../../include/open-logic-section]{subfiles}

\begin{document}

\olfileid{cmp}{rec}{gen}
\olsection{સામાન્ય પુનરાવર્તી વિધેયો}

સંપૂર્ણ વિધેયોનો ગણ મેળવવાની બીજી રીત પણ છે. સંપૂર્ણ વિધેય
$f(x,\vec z)$ને \emph{નિયમિત} કહીએ, જો પ્રાકૃતિક સંખ્યાઓના દરેક
અનુક્રમ $\vec z$ માટે એવો $x$ હોય કે $f(x,\vec z) = 0$ થાય.
બીજા શબ્દોમાં, નિયમિત વિધેયો ચોક્કસ એવાં વિધેયો છે જેના પર
અનિયત મર્યાદા સુધીની શોધ લાગુ કરતાં સંપૂર્ણ વિધેય મળે છે.
સાવધાનીપૂર્વક, અનિયત મર્યાદા સુધીની શોધને નિયમિત વિધેયો પૂરતી
મર્યાદિત કરી શકીએ:

\begin{defn}
\ollabel{defn:general-recursive}
\emph{સામાન્ય પુનરાવર્તી વિધેયો}નો ગણ એટલે પ્રાકૃતિક
સંખ્યાઓમાંથી પ્રાકૃતિક સંખ્યાઓમાં જતાં (જુદી જુદી
આર્ગ્યુમેન્ટ-સંખ્યા ધરાવતાં) વિધેયોનો સૌથી નાનો એવો ગણ, જેમાં
શૂન્ય વિધેય, ઉત્તરગામી વિધેય અને પ્રક્ષેપો છે તથા જે સંયોજન,
આદિમ પુનરાવર્તન અને \emph{નિયમિત} વિધેયો પર લાગુ પડતી અનિયત
મર્યાદા સુધીની શોધ હેઠળ બંધ છે.
\end{defn}

સ્પષ્ટ છે કે દરેક સામાન્ય પુનરાવર્તી વિધેય સંપૂર્ણ છે.
\olref{defn:general-recursive} અને
\olref[par]{defn:recursive-fn} વચ્ચેનો તફાવત એ છે કે પછીની
વ્યાખ્યામાં વચ્ચેના પગલાઓમાં આંશિક પુનરાવર્તી વિધેયો વાપરવાની
છૂટ છે; અંતે મળતું વિધેય સંપૂર્ણ હોય એટલી જ શરત છે. તેથી
``સામાન્ય'' શબ્દ ઐતિહાસિક અવશેષ છે અને નામ તરીકે ભ્રામક છે:
ઉપરથી જોતાં \olref{defn:general-recursive} એ
\olref[par]{defn:recursive-fn} કરતાં \emph{ઓછી} સામાન્ય લાગે છે.
પરંતુ સદનસીબે આ તફાવત માત્ર દેખીતો છે: વ્યાખ્યાઓ જુદી હોવા
છતાં સામાન્ય પુનરાવર્તી વિધેયોનો ગણ અને પુનરાવર્તી વિધેયોનો
ગણ એક જ છે.

\end{document}
